\documentclass[11pt]{article}
\usepackage[margin=1in]{geometry}
\usepackage{amsmath}
\usepackage[hidelinks]{hyperref}

\title{Chromatic polynomials do not determine configuration Euler characteristics\\[2pt]
  \large Disproof of Conjecture 00000004093}
\author{}
\date{}

\begin{document}
\maketitle

\noindent
The four-vertex path $P_4$ and the three-leaf star $K_{1,3}$ have the same
chromatic polynomial. Their two-particle configuration spaces have
different ordinary Euler characteristics, both for labeled particles and
for unlabeled particles.

\section*{The graphs and their chromatic polynomials}

On vertices $0,1,2,3$, the path has edges $01,12,23$ and the star has edges
$01,12,13$. In either tree, a root has $k$ possible colors and each of the
three remaining vertices has $k-1$ choices once its parent has been colored.
Thus both graphs have chromatic polynomial $P(t)=t(t-1)^3$. In particular,
both have $P(-2)=54$ and $P(-1)=8$.

\section*{The configuration spaces}

For a graph $G$ with geometric realization $|G|$, let
$F_2(G)=\{(x,y)\in |G|^2:x\ne y\}$ and let $U_2(G)$ be its quotient by the
coordinate swap. Write $D_2(G)$ for the finite cubical subcomplex with cells
$c\times d$ whose closures in $|G|$ are disjoint. Here $c$ and $d$ are
vertices or open edges of $G$.

The stable equivalence theorem of Prue and Scrimshaw~\cite[Theorem~3.2]{ps}
gives a deformation retraction from $F_n(G)$ to its discretization for a
finite connected graph with at least $n$ vertices, provided that every path
between distinct vertices of degree other than two has at least $n-1$
edges and every noncontractible closed path has at least $n+1$ edges.
The theorem also gives the corresponding retraction for unlabeled particles.
Both graphs satisfy these conditions at $n=2$: a path between distinct
vertices has at least one edge, and a tree has no noncontractible loops.

For these two graphs the retraction can also be seen directly. On the
product of an edge with itself, continuously move the lesser and greater
positions to the two endpoints by linear interpolation. On the product of
two incident distinct edges, let $u,v\in[0,1]$ be coordinates measured from
their common endpoint and interpolate linearly to
$(u/\max(u,v),v/\max(u,v))$. The maximum is positive in a configuration.
On the product of disjoint edges keep the configuration fixed. These
motions preserve distinctness of the particles. When one particle is at
a vertex, the formulas agree on intersecting products of edges. A vertex
incident to the other particle's edge stays fixed and the other particle
moves toward its far endpoint. A nonincident vertex gives a stationary
configuration. Pairs of distinct vertices also stay fixed. The local
motions therefore glue to a continuous deformation, end in $D_2(G)$, and
fix it pointwise. They commute with the coordinate swap, so they descend
to the unlabeled space.

\section*{Cell counts and the contradiction}

For either graph there are $4\cdot3=12$ ordered vertex pairs. Each edge
has two vertices outside its closure, so there are $2\cdot3\cdot2=12$
one-cells, counting both placements of the vertex coordinate. The path
has exactly one pair of disjoint edges, $01$ and $23$, giving two ordered
two-cells. Every pair of star edges meets at vertex $1$, so the star has
no two-cells. Coordinate swap pairs distinct cells of the same dimension.
The unordered quotient has one cell per pair, giving the counts below.
\[
\begin{array}{c|rrr|r}
\text{Space} & \text{0-cells} & \text{1-cells} & \text{2-cells} & \chi \\ \hline
F_2(P_4)     & 12 & 12 & 2 & 2 \\
F_2(K_{1,3}) & 12 & 12 & 0 & 0 \\
U_2(P_4)     &  6 &  6 & 1 & 1 \\
U_2(K_{1,3}) &  6 &  6 & 0 & 0
\end{array}
\]
These are counts in the respective deformation retracts. The factor of two
between the ordered and unordered Euler characteristics comes from the
two-sheeted covering. A universal expression in $P_G(-n)$ and $P_G(-1)$
at $n=2$ would give the same value for both graphs, contradicting either
pair of Euler characteristics in the table.

\section*{Lean verification}

The Lean 4 project enumerates products of cells with disjoint closures
and one representative of each swap orbit. It checks the counts above.
It also gives an explicit bijection between proper colorings of the two
graphs for every finite palette. Keep the first three colors fixed and
apply the transposition of the colors on vertices $1$ and $2$ to the color
on vertex $3$. The same transformation recovers the original coloring.
The theorem \texttt{C4093.no\_chromatic\_rule} rules out any recovery of
either cell Euler characteristic from this chromatic data. Run
\texttt{lake build} in the \texttt{lean} directory to check it.

\begin{thebibliography}{1}
\bibitem{ps}
P.~Prue and T.~Scrimshaw, \emph{Abrams's stable equivalence for graph braid groups}.
\href{https://arxiv.org/abs/0909.5511}{arXiv:0909.5511}.
\end{thebibliography}

\end{document}
